\begin{table}[H]\centering
\caption{\textbf{Descriptive ranking: the weekend dip across all 21 two-day placebos}}
\label{tab:permutation}\small
\small\setlength{\tabcolsep}{4pt}\resizebox{\ifdim\width>\linewidth \linewidth\else\width\fi}{!}{%
\begin{tabular}{lc}
\toprule
Pseudo rest-day pair & Private cesarean dip (pp) \\
\midrule
Sun+Sat (true weekend) & -8.11 \\
Sun+Fri & -3.90 \\
Sun+Thu & -3.62 \\
Sun+Wed & -3.20 \\
Sun+Tue & -3.16 \\
Sun+Mon & -1.78 \\
Fri+Sat & -1.58 \\
Thu+Sat & -1.36 \\
Wed+Sat & -0.98 \\
Tue+Sat & -0.95 \\
Mon+Sat & 0.26 \\
Thu+Fri & 1.37 \\
Wed+Fri & 1.70 \\
Tue+Fri & 1.72 \\
Wed+Thu & 1.87 \\
Tue+Thu & 1.90 \\
Tue+Wed & 2.21 \\
Mon+Fri & 2.77 \\
Mon+Thu & 2.93 \\
Mon+Wed & 3.24 \\
Mon+Tue & 3.27 \\
\bottomrule
\end{tabular}
}
\\[2pt]\footnotesize\textit{Notes:} Each row re-estimates the private-sector weekend term of Equation~\eqref{eq:scheduling}, treating a different pair of weekdays as the ``rest days'' (SINASC 2010--2024, municipality-date cells, municipality and year fixed effects, weighted by births). The true weekend (Sun+Sat) is the most negative of all 21 placebos (rank 1 of 21). Because the days of the week are not exchangeable under a known assignment mechanism, this is a descriptive ranking, not an exact randomization-inference $p$-value.
\end{table}
